\section{本章纲要：强国的逻辑与帝国的难题}

从\eventref{WQ-0453-ZHI}{三家分晋}到\eventref{QIN-0209-FALL}{秦亡}，战国秦的长线不是“谁最会打仗”，而是谁最能把社会资源持续转化为国家行动。晋阳城下改了水道，洛邑朝会改了名分，县廷与军营则一点点改了人、粮和命令的去向。魏、楚、韩、秦等国的改革共同抬高了竞争门槛；秦把这套能力结合得最彻底，因而完成统一；它又在和平时期以过高强度使用这套能力，因而迅速失去承受者。

这一章不把七国排成一张胜负榜。魏的先发、楚的纵深、齐的富庶、赵与燕的北方边线、韩的通道困局，以及残余周廷所余的名分，各自规定了改革、结盟和用兵的成本。合纵可以在关外集结，却未必能共给一条粮道；一座城可以攻下，却未必能把新居民、旧官吏和田租接入胜者的簿书。秦的优势由关中、巴蜀、县邑、军功与连续战争共同织成，不是任何一条法令或一位将领单独造成。

读者将沿共享年表看三家分晋、称王、伊阙、长平与周室终局如何相连；再回到各国 dossier，看每次选择落在何处山口、都城、河道与继承关系上；最后在统一与秦亡中检验同一套征发能力为何既能合天下、又会耗尽地方承受力。传世史书提供年代与人物线索，简牍、器物、战场和历史地理则不断提醒我们：后人保存下来的名言、兵数和谋略，并不是现场本身。

下一章的汉，不会简单“推翻秦的一切”。它会继承郡县、法律、道路和大一统的想象，同时寻找一种比秦更能长期维持的统治节奏。
